\section{Proof of the main theorem}\label{sec:proofmain}

Theorem~\ref{thm:main} follows from five results proved above: the variance bound
(Theorem~\ref{thm:variance}), the integer lower bound (Theorem~\ref{thm:integer}), the tree
bound for the fractional optimum (Theorem~\ref{thm:lp}), the tree process
(Theorem~\ref{thm:process}) and the existence and cost of a fixed point
(Theorem~\ref{thm:existence}). The section ends with a consequence of
Theorems~\ref{thm:integer} and~\ref{thm:lp} for the tree itself.

\subsection{The derivation}\label{sec:proofmain-derivation}

\begin{proof}[Proof of Theorem~\ref{thm:main}]
\emph{The fractional deficit.} By Theorem~\ref{thm:existence} the map $\RDE$ has a fixed
point $\mu$, and $9\,\E X\le\alphalp$ for $X\sim\mu$. By Theorem~\ref{thm:process} there is
an invariant perfect matching $m$ on $\tree$ with cost $c(m)=9\,\E X\le\alphalp$. By
Theorem~\ref{thm:lp},
\[
\E\PhiLP\le(\alphalp+o(1))\,q^{1/3}.
\]

\emph{The integer deficit.} By Theorem~\ref{thm:integer},
$\P(\Phiq\ge\betaint q^{1/3})\to1$. Since $\Phiq\ge0$,
\[
\E\Phiq\ \ge\ \betaint\,q^{1/3}\,\P\bigl(\Phiq\ge\betaint q^{1/3}\bigr)
=(\betaint-o(1))\,q^{1/3}.
\]

\emph{Part \textup{(i)}.} By \eqref{eq:model-gap}, $\Dq=\Phiq-\PhiLP$, so
\[
\E\Dq=\E\Phiq-\E\PhiLP\ \ge\ (\betaint-\alphalp-o(1))\,q^{1/3},
\]
and $\betaint-\alphalp=2.302-2.2519=0.0501$. In particular $\E\Dq\to\infty$.

\emph{Part \textup{(ii)}.} By Theorem~\ref{thm:variance}(ii) with $\eps=1/12$, $\Var\Dq$ is
bounded for large $q$ (it is $O(q^{-1/12})$). Fix $t>0$. By Chebyshev's inequality and
part~(i),
\[
\P\bigl(|\Dq-\E\Dq|>t\,\E\Dq\bigr)\le\frac{\Var\Dq}{t^2(\E\Dq)^2}\longrightarrow0
\qquad(q\to\infty).
\]
This holds for every fixed $t>0$, which is the statement $\Dq/\E\Dq\to1$ in probability.

\emph{Part \textup{(iii)}.} By Proposition~\ref{prop:model-core} the core is nonempty
exactly when $\Dq=0$. For large $q$ we have $\E\Dq>0$, and then the event $\{\Dq=0\}$ is
contained in $\{|\Dq-\E\Dq|>\tfrac12\E\Dq\}$. Take $t=1/2$ in part~(ii).
\end{proof}

The argument gives $\P(\Dq=0)\le4\Var\Dq/(\E\Dq)^2=O(q^{-5/6+\eps})$ for every $\eps>0$, by
Theorem~\ref{thm:variance}(ii) and part~(i); compare Corollary~\ref{cor:variance-core}. The
implied constants and thresholds are those of Theorems~\ref{thm:variance},
\ref{thm:integer} and~\ref{thm:lp}, and none of them is effective at any size that can be
simulated.

\subsection{The matching is not a function of the tree}\label{sec:proofmain-endogeny}

The proof of Theorem~\ref{thm:lp} transfers to the market only the conditional expectation
of the matching given the tree and its costs, and for that reason it bounds $\PhiLP$ and
not $\Phiq$. If the matching were a function of the tree, the same construction would
produce assignments, and Theorem~\ref{thm:integer} forbids this at cost below $\betaint$.

\begin{corollary}\label{cor:proofmain-endogeny}
Let $m$ be an invariant perfect matching on $\tree$ with $c(m)<\betaint$. Then $m$ is not a
measurable function of the marked tree: with $\ell_1,\ell_2,\dots$ the teams at the root in
increasing order of cost, there are no measurable functions $h_j$ with values in $\{0,1\}$
such that $m_{\ell_j}=h_j(\tree,o,\xi)$ almost surely for every $j$. In particular the
invariant perfect matching of Theorem~\ref{thm:process} built from a fixed point of $\RDE$,
whose cost is at most $\alphalp<\betaint$ by Theorem~\ref{thm:existence}, uses randomness
beyond the tree and its costs.
\end{corollary}

\begin{proof}
Suppose that $m_{\ell_j}=h_j(\tree,o,\xi)$ almost surely with measurable $h_j$ with values in
$\{0,1\}$, and that $c(m)<\infty$. We show $c(m)\ge\betaint$. We use the notation of
Section~\ref{sec:lp}; the functions $h_j$ may be taken as the functions of
Section~\ref{sec:lp-conditional}, since $\E[m_{\ell_j}\mid\tree,o,\xi]=h_j$.

\emph{A thresholded rule.} Let $g_r$ be the rule of Lemma~\ref{lem:lp-condrule} and put
\[
\tilde g_{r,j}=\ind\{g_{r,j}>\tfrac12\},\qquad j\ge1 .
\]
Since $\sum_jg_{r,j}\le1$, at most one $\tilde g_{r,j}(B)$ equals $1$ for each ball $B$, and
$\tilde g_{r,j}$ vanishes where $g_{r,j}$ does. So $\tilde g_r$ is a rule of range $r$ at
level $r$ in the sense of Definition~\ref{def:lp-rule}, with values in $\{0,1\}$, and the
weights $x^{\tilde g_r}$ of \eqref{eq:lp-minpacking} take values in $\{0,1\}$: a team is
kept when its conditional probability exceeds $1/2$ from the point of view of each of its
three agents.

\emph{On the tree.} Write $G_{r,j}=g_{r,j}(B_r(\tree_r,o))$ and
$\tilde G_{r,j}=\ind\{G_{r,j}>\tfrac12\}$. Almost surely exactly one $h_j$ equals $1$; let
$J$ be its index, so that $\cost_{\ell_J}$ is the cost of the team of the root in $m$ and
$\E\,\cost_{\ell_J}=c(m)$. By Lemma~\ref{lem:lp-condrule}(b), $G_{r,J}\to h_J=1$ almost
surely. So almost surely, for all large $r$, $\tilde G_{r,J}=1$, hence $\tilde G_{r,j}=0$
for $j\ne J$, that is, $\tilde G_{r,j}=h_j$ for every $j$. Two consequences follow.

First, by Lemma~\ref{lem:lp-tree},
\[
\tilde\eps_r:=\E\,S^{\tilde g_r}\bigl(B_{r+1}(\tree_r,o)\bigr)
\le3\,\E\sum_{j\ge1}\bigl(h_j-\tilde G_{r,j}\bigr)_+\longrightarrow0,
\]
by bounded convergence: the sum is at most $1$ and vanishes for large $r$.

Secondly, by Lemma~\ref{lem:lp-tree},
$\tilde c_r:=\E\,C^{\tilde g_r}(B_{r+1}(\tree_r,o))\le\E\,Y_r$ with
$Y_r=\sum_j\cost_{\ell_j}\tilde G_{r,j}$, and $Y_r=\cost_{\ell_J}$ for all large $r$, almost
surely. Moreover $\tilde G_{r,j}\le2G_{r,j}$, and by \eqref{eq:lp-condrule}, since
$\cost_{\ell_j}\ind\{\cost_{\ell_j}\le r\}$ is nonnegative and $\mathcal F_r$-measurable,
\[
Y_r\le2\sum_{j\ge1}\cost_{\ell_j}G_{r,j}
=2\,\E\Bigl[\sum_{j\ge1}\cost_{\ell_j}\ind\{\cost_{\ell_j}\le r\}\,m_{\ell_j}\Bigm|
\mathcal F_r\Bigr]\le2\,\E\bigl[\cost_{\ell_J}\bigm|\mathcal F_r\bigr].
\]
The conditional expectations of the integrable variable $\cost_{\ell_J}$ form a uniformly
integrable family. Hence $(Y_r)_r$ is uniformly integrable, and
$\E\,Y_r\to\E\,\cost_{\ell_J}=c(m)$. So $\limsup_r\tilde c_r\le c(m)$.

\emph{In the market.} Fix $r\ge1$, $\eta\in(0,1)$ and $q\ge r^{3/2}$, and repeat the proof
of Theorem~\ref{thm:lp} with the rule $\tilde g_r$ in place of $g_r$. The vector
$x=x^{\tilde g_r}$, computed in the market $u$, takes values in $\{0,1\}$, so the last
statement of Lemma~\ref{lem:lp-patch} applies, and with Lemma~\ref{lem:lp-market} it gives
\[
\frac{\E\Phiq}{q^{1/3}}\le\frac{\tilde c_r+r\,\eps_1(q,r,r+1)}{1-\eta}
+\frac{C_0}{\eta}\bigl(\tilde\eps_r+\eps_1(q,r,r+1)\bigr)^{1/3}.
\]
Let $q\to\infty$, then $r\to\infty$, then $\eta\to0$:
$\limsup_q\E\Phiq/q^{1/3}\le c(m)$.

\emph{Conclusion.} By Theorem~\ref{thm:integer} and $\Phiq\ge0$, as in the proof of
Theorem~\ref{thm:main}, $\liminf_q\E\Phiq/q^{1/3}\ge\betaint$. Hence $c(m)\ge\betaint$.
\end{proof}

\begin{remark}\label{rem:proofmain-endogeny}
The matching of Theorem~\ref{thm:process} is a function of the tree, the costs and a family
of messages that solve the recursion on the tree (Remark~\ref{rem:tree-factor}). By
Corollary~\ref{cor:proofmain-endogeny} the messages cannot all be measurable functions of
the marked tree: in the usual terminology, the recursive tree process attached to a fixed
point of $\RDE$ is not endogenous. No statement in this paper relies on endogeny.
\end{remark}
